\Article{14}
Thus, when the function \(y\) has been reduced to the form of a sum of terms such as
\[
Ax^\alpha x'{}^\beta x''{}^\gamma\mathord{.}\,\mathord{.}\,\mathord{.}\,\mathord{.}
\]
the value of the integral \(\int y\psi y.\mathrm{d}y\) extended over all values of \(y\), or the mean value of \(y\), will be equal to the sum of terms
\[
A\times\int x^\alpha\varphi x.\mathrm{d}x\times\int x'{}^\beta\varphi x'.\mathrm{d}x'\times\int x''{}^\gamma\varphi x''.\mathrm{d}x''\mathord{.}\,\mathord{.}\,\mathord{.}\,\mathord{.}
\]
where the integrations must be extended from \(x=-\infty\) to \(x=+\infty\), from \(x'=-\infty\) to \(x'=+\infty\) etc.; or, what amounts to the same thing, to the sum of the terms that arise when the individual powers \(x^\alpha,x'{}^\beta,x''{}^\gamma\) etc. are replaced by their mean values. The truth of this very important theorem could also easily have been derived from other considerations.

\Article{15}
Let us apply what we set out in the preceding article to the special case in which
\[
y=\frac{xx+x'x'+x''x''+\text{etc.}}{\sigma}
\]
where \(\sigma\) denotes the number of terms in the numerator. The mean value of \(y\) is immediately found here \(=mm\), taking the symbol \(m\) in the same sense as above. The actual value of \(y\) in a definite case may indeed turn out greater or less than the mean, just as may the actual value of the single term \(xx\): but the probability that the random value of \(y\) will not differ appreciably from the mean \(mm\) will approach certainty more and more as the number \(\sigma\) increases. To make this clearer, since it is not in our power to determine the probability itself exactly, let us investigate the mean error to be feared when we suppose \(y=mm\). Evidently, by the principles of art. 6, this error will be the square root of the mean value of the function
\[
\left(\frac{xx+x'x'+x''x''+\text{etc.}}{\sigma}-mm\right)^2
\]
to find which it is enough to observe that the mean value of a term such as \(\frac{x^4}{\sigma\sigma}\) is \(=\frac{n^4}{\sigma\sigma}\) (using the symbol \(n\) in the sense of art. 11), while the mean value of a term \SourceLeaf{020}such as \(\frac{2xxx'x'}{\sigma\sigma}\) becomes \(=\frac{2m^4}{\sigma\sigma}\), whence the mean value of that function is very easily deduced
\[
=\frac{n^4-m^4}{\sigma}
\]

From this we learn that, if a sufficiently large supply of mutually independent random errors \(x,x',x''\) etc. is available, an approximate value of \(m\) can be obtained from them with great certainty by the formula
\[
m=\BareRoot{\frac{(xx+x'x'+x''x''+\text{etc.})}{\sigma}}
\]
and that the mean error to be feared in this determination, with regard to the square \(mm\), is
\[
=\BareRoot{\frac{n^4-m^4}{\sigma}}
\]
Moreover, since the latter formula involves the quantity \(n\), if the only aim is to form some idea of the degree of precision of this determination, it will be enough to adopt some hypothesis concerning the function \(\varphi\). For example, under the third hypothesis of artt. 9, 11, that error becomes \(=mm\BareRoot{\frac2\sigma}\). If this is less appealing, an approximate value of \(n^4\) can be obtained from the errors themselves with the aid of the formula
\[
\frac{x^4+x'{}^4+x''{}^4+\text{etc.}}{\sigma}
\]
In general, however, we can assert that doubling the precision in this determination requires a fourfold supply of errors, or that the weight of the determination is proportional to the number \(\sigma\).

In exactly the same way, if the errors of observations involve a constant part, an approximate value of this part can be inferred from the arithmetic mean of many errors with greater confidence the greater their number. And the mean error to be feared in this determination will be expressed by
\[
\BareRoot{\frac{mm-kk}{\sigma}}
\]
if \(k\) denotes the constant part itself and \(m\) the mean error of the observations not yet freed of the constant part, or simply by \(\frac{m}{\BareRoot{\sigma}}\) if \(m\) denotes the mean error of the observations freed of the constant part (see art. 8).

\clearpage
